
\documentclass{article}
\usepackage{colortbl}
\usepackage{makecell}
\usepackage{multirow}
\usepackage{supertabular}

\begin{document}

\newcounter{utterance}

\twocolumn

{ \footnotesize  \setcounter{utterance}{1}
\setlength{\tabcolsep}{0pt}
\begin{supertabular}{c@{$\;$}|p{.15\linewidth}@{}p{.15\linewidth}p{.15\linewidth}p{.15\linewidth}p{.15\linewidth}p{.15\linewidth}}

    \# & $\;$A & \multicolumn{4}{c}{Game Master} & $\;\:$B\\
    \hline 

    \theutterance \stepcounter{utterance}  

    & & \multicolumn{4}{p{0.6\linewidth}}{\cellcolor[rgb]{0.9,0.9,0.9}{%
	\makecell[{{p{\linewidth}}}]{% 
	  \tt {\tiny [A$\langle$GM]}  
	 Du nimmst an einem kollaborativen Verhandlungspiel Teil.\\ \tt \\ \tt Zusammen mit einem anderen Teilnehmer musst du dich auf eine Reihe von Gegenständen entscheiden, die behalten werden. Jeder von euch hat eine persönliche Verteilung über die Wichtigkeit der einzelnen Gegenstände. Jeder von euch hat eine eigene Meinung darüber, wie wichtig jeder einzelne Gegenstand ist (Gegenstandswichtigkeit). Du kennst die Wichtigkeitsverteilung des anderen Spielers nicht. Zusätzlich siehst du, wie viel Aufwand jeder Gegenstand verursacht.  \\ \tt Ihr dürft euch nur auf eine Reihe von Gegenständen einigen, wenn der Gesamtaufwand der ausgewählten Gegenstände den Maximalaufwand nicht überschreitet:\\ \tt \\ \tt Maximalaufwand = 3358\\ \tt \\ \tt Hier sind die einzelnen Aufwände der Gegenstände:\\ \tt \\ \tt Aufwand der Gegenstände = {"A07": 171, "A36": 933, "B73": 552, "C45": 219, "B93": 275, "B98": 778, "A87": 341, "A63": 615, "B05": 519, "C13": 862, "B16": 262, "B33": 377, "C34": 347, "C89": 349, "C11": 117}\\ \tt \\ \tt Hier ist deine persönliche Verteilung der Wichtigkeit der einzelnen Gegenstände:\\ \tt \\ \tt Werte der Gegenstandswichtigkeit = {"A07": 138, "A36": 583, "B73": 868, "C45": 822, "B93": 783, "B98": 65, "A87": 262, "A63": 121, "B05": 508, "C13": 780, "B16": 461, "B33": 484, "C34": 668, "C89": 389, "C11": 808}\\ \tt \\ \tt Ziel:\\ \tt \\ \tt Dein Ziel ist es, eine Reihe von Gegenständen auszuhandeln, die dir möglichst viel bringt (d. h. Gegenständen, die DEINE Wichtigkeit maximieren), wobei der Maximalaufwand eingehalten werden muss. Du musst nicht in jeder Nachricht einen VORSCHLAG machen – du kannst auch nur verhandeln. Alle Taktiken sind erlaubt!\\ \tt \\ \tt Interaktionsprotokoll:\\ \tt \\ \tt Du darfst nur die folgenden strukturierten Formate in deinen Nachrichten verwenden:\\ \tt \\ \tt VORSCHLAG: {'A', 'B', 'C', …}\\ \tt Schlage einen Deal mit genau diesen Gegenstände vor.\\ \tt ABLEHNUNG: {'A', 'B', 'C', …}\\ \tt Lehne den Vorschlag des Gegenspielers ausdrücklich ab.\\ \tt ARGUMENT: {'...'}\\ \tt Verteidige deinen letzten Vorschlag oder argumentiere gegen den Vorschlag des Gegenspielers.\\ \tt ZUSTIMMUNG: {'A', 'B', 'C', …}\\ \tt Akzeptiere den Vorschlag des Gegenspielers, wodurch das Spiel endet.\\ \tt STRATEGISCHE ÜBERLEGUNGEN: {'...'}\\ \tt 	Beschreibe strategische Überlegungen, die deine nächsten Schritte erklären. Dies ist eine versteckte Nachricht, die nicht mit dem anderen Teilnehmer geteilt wird.\\ \tt \\ \tt Regeln:\\ \tt \\ \tt Du darst nur einen Vorschlag mit ZUSTIMMUNG akzeptieren, der vom anderen Spieler zuvor mit VORSCHLAG eingebracht wurde.\\ \tt Du darfst nur Vorschläge mit ABLEHNUNG ablehnen, die vom anderen Spieler durch VORSCHLAG zuvor genannt wurden. \\ \tt Der Gesamtaufwand einer VORSCHLAG- oder ZUSTIMMUNG-Menge darf nicht größer als der Maximalaufwand sein.  \\ \tt Offenbare deine versteckte Wichtigkeitsverteilung nicht.\\ \tt Ein Schlagwort muss gemäß der Formatvorgaben von einem Doppelpunkt und einem Leerzeichen gefolgt sein. Das Argument muss eine Python-Menge mit 0 oder mehr Strings sein.  \\ \tt Die Form muss also SCHLAGWORT: {...} sein.\\ \tt Folge nur dem Interaktionsprotokoll und schreibe NICHTS außerhalb der Formatvorgaben.\\ \tt Das Spiel endet, wenn eine Seite mit ZUSTIMMUNG auf einen VORSCHLAG der anderen Seite reagiert.  \\ \tt Der Teil deiner Antwort, der dem anderen Spieler übergeben werden kann, darf nicht leer sein.  \\ \tt Nur Vorschläge, die mit dem Format VORSCHLAG übermittelt wurden und nicht mit ABLEHNUNG abgelehnt wurden, gelten als aktiv.  \\ \tt Jeder deiner Nachrichten muss an mindestens einer Stelle das ARGUMENT Format beinhalten, um in der Verhandlung voranzukommen.\\ \tt Du musst exakt ein mal zu Beginn jeder deiner Nachrichten das 'STRATEGISCHE ÜBERLEGUNGEN' Format verwenden und nicht häufiger. Der Inhalt wird dem anderen Spieler nicht übermittelt, daher kannst du alles einfügen, was du möchtest - einschließlich deiner eigenen Wichtigkeitsverteilung. Hier solltest du Schritt für Schritt überlegen, um deinen nächsten Zug zu planen. 
	  } 
	   } 
	   } 
	 & \\ 
 

    \theutterance \stepcounter{utterance}  

    & \multicolumn{4}{p{0.6\linewidth}}{\cellcolor[rgb]{0.8,1,0.9}{%
	\makecell[{{p{\linewidth}}}]{% 
	  \tt {\tiny [A$\rangle$GM]}  
	 STRATEGISCHE ÜBERLEGUNGEN: {'Ich muss eine Sammlung von Gegenständen auswählen, die meine Gesamtwichtigkeit maximiert und gleichzeitig den Maximalaufwand von 3358 nicht überschreitet. Meine wichtigsten Gegenstände sind B73 (868), C11 (808), C45 (822) und B93 (783) mit moderaten Aufwänden. Allerdings habe ich Gegenstände mit niedrigen Werten wie B98 (65) und A63 (121), die ich opfern kann. Da ich die Präferenzen meines Partners nicht kenne, sollte ich einen initialen Vorschlag machen, der meine hochwertigen Gegenstände enthält, aber Raum für Verhandlungen lässt.'}\\ \tt \\ \tt VORSCHLAG: {'A07', 'A36', 'B73', 'C45', 'B93', 'A87', 'B05', 'C13', 'B16', 'B33', 'C34', 'C89', 'C11'}\\ \tt \\ \tt ARGUMENT: {'Ich schlage eine Auswahl vor, die eine breite Palette von Gegenständen umfasst und hohen Wert für mich bietet. Dieser Vorschlag bleibt unter dem Maximalaufwand und bietet Flexibilität für Anpassungen basierend auf Ihren Prioritäten.'} 
	  } 
	   } 
	   } 
	 & & \\ 
 

    \theutterance \stepcounter{utterance}  

    & & \multicolumn{4}{p{0.6\linewidth}}{\cellcolor[rgb]{0.9,0.9,0.9}{%
	\makecell[{{p{\linewidth}}}]{% 
	  \tt {\tiny [GM$\rangle$B]}  
	 Du nimmst an einem kollaborativen Verhandlungspiel Teil.\\ \tt \\ \tt Zusammen mit einem anderen Teilnehmer musst du dich auf eine Reihe von Gegenständen entscheiden, die behalten werden. Jeder von euch hat eine persönliche Verteilung über die Wichtigkeit der einzelnen Gegenstände. Jeder von euch hat eine eigene Meinung darüber, wie wichtig jeder einzelne Gegenstand ist (Gegenstandswichtigkeit). Du kennst die Wichtigkeitsverteilung des anderen Spielers nicht. Zusätzlich siehst du, wie viel Aufwand jeder Gegenstand verursacht.  \\ \tt Ihr dürft euch nur auf eine Reihe von Gegenständen einigen, wenn der Gesamtaufwand der ausgewählten Gegenstände den Maximalaufwand nicht überschreitet:\\ \tt \\ \tt Maximalaufwand = 3358\\ \tt \\ \tt Hier sind die einzelnen Aufwände der Gegenstände:\\ \tt \\ \tt Aufwand der Gegenstände = {"A07": 171, "A36": 933, "B73": 552, "C45": 219, "B93": 275, "B98": 778, "A87": 341, "A63": 615, "B05": 519, "C13": 862, "B16": 262, "B33": 377, "C34": 347, "C89": 349, "C11": 117}\\ \tt \\ \tt Hier ist deine persönliche Verteilung der Wichtigkeit der einzelnen Gegenstände:\\ \tt \\ \tt Werte der Gegenstandswichtigkeit = {"A07": 138, "A36": 583, "B73": 868, "C45": 822, "B93": 783, "B98": 65, "A87": 262, "A63": 121, "B05": 508, "C13": 780, "B16": 461, "B33": 484, "C34": 668, "C89": 389, "C11": 808}\\ \tt \\ \tt Ziel:\\ \tt \\ \tt Dein Ziel ist es, eine Reihe von Gegenständen auszuhandeln, die dir möglichst viel bringt (d. h. Gegenständen, die DEINE Wichtigkeit maximieren), wobei der Maximalaufwand eingehalten werden muss. Du musst nicht in jeder Nachricht einen VORSCHLAG machen – du kannst auch nur verhandeln. Alle Taktiken sind erlaubt!\\ \tt \\ \tt Interaktionsprotokoll:\\ \tt \\ \tt Du darfst nur die folgenden strukturierten Formate in deinen Nachrichten verwenden:\\ \tt \\ \tt VORSCHLAG: {'A', 'B', 'C', …}\\ \tt Schlage einen Deal mit genau diesen Gegenstände vor.\\ \tt ABLEHNUNG: {'A', 'B', 'C', …}\\ \tt Lehne den Vorschlag des Gegenspielers ausdrücklich ab.\\ \tt ARGUMENT: {'...'}\\ \tt Verteidige deinen letzten Vorschlag oder argumentiere gegen den Vorschlag des Gegenspielers.\\ \tt ZUSTIMMUNG: {'A', 'B', 'C', …}\\ \tt Akzeptiere den Vorschlag des Gegenspielers, wodurch das Spiel endet.\\ \tt STRATEGISCHE ÜBERLEGUNGEN: {'...'}\\ \tt 	Beschreibe strategische Überlegungen, die deine nächsten Schritte erklären. Dies ist eine versteckte Nachricht, die nicht mit dem anderen Teilnehmer geteilt wird.\\ \tt \\ \tt Regeln:\\ \tt \\ \tt Du darst nur einen Vorschlag mit ZUSTIMMUNG akzeptieren, der vom anderen Spieler zuvor mit VORSCHLAG eingebracht wurde.\\ \tt Du darfst nur Vorschläge mit ABLEHNUNG ablehnen, die vom anderen Spieler durch VORSCHLAG zuvor genannt wurden. \\ \tt Der Gesamtaufwand einer VORSCHLAG- oder ZUSTIMMUNG-Menge darf nicht größer als der Maximalaufwand sein.  \\ \tt Offenbare deine versteckte Wichtigkeitsverteilung nicht.\\ \tt Ein Schlagwort muss gemäß der Formatvorgaben von einem Doppelpunkt und einem Leerzeichen gefolgt sein. Das Argument muss eine Python-Menge mit 0 oder mehr Strings sein.  \\ \tt Die Form muss also SCHLAGWORT: {...} sein.\\ \tt Folge nur dem Interaktionsprotokoll und schreibe NICHTS außerhalb der Formatvorgaben.\\ \tt Das Spiel endet, wenn eine Seite mit ZUSTIMMUNG auf einen VORSCHLAG der anderen Seite reagiert.  \\ \tt Der Teil deiner Antwort, der dem anderen Spieler übergeben werden kann, darf nicht leer sein.  \\ \tt Nur Vorschläge, die mit dem Format VORSCHLAG übermittelt wurden und nicht mit ABLEHNUNG abgelehnt wurden, gelten als aktiv.  \\ \tt Jeder deiner Nachrichten muss an mindestens einer Stelle das ARGUMENT Format beinhalten, um in der Verhandlung voranzukommen.\\ \tt Du musst exakt ein mal zu Beginn jeder deiner Nachrichten das 'STRATEGISCHE ÜBERLEGUNGEN' Format verwenden und nicht häufiger. Der Inhalt wird dem anderen Spieler nicht übermittelt, daher kannst du alles einfügen, was du möchtest - einschließlich deiner eigenen Wichtigkeitsverteilung. Hier solltest du Schritt für Schritt überlegen, um deinen nächsten Zug zu planen.\\ \tt \\ \tt Du erhältst nun die erste Nachricht des anderen Spielers.\\ \tt \\ \tt VORSCHLAG: {'A07', 'A36', 'B73', 'C45', 'B93', 'A87', 'B05', 'C13', 'B16', 'B33', 'C34', 'C89', 'C11'}\\ \tt \\ \tt ARGUMENT: {'Ich schlage eine Auswahl vor, die eine breite Palette von Gegenständen umfasst und hohen Wert für mich bietet. Dieser Vorschlag bleibt unter dem Maximalaufwand und bietet Flexibilität für Anpassungen basierend auf Ihren Prioritäten.'} 
	  } 
	   } 
	   } 
	 & \\ 
 

    \theutterance \stepcounter{utterance}  

    & & & \multicolumn{4}{p{0.6\linewidth}}{\cellcolor[rgb]{1,0.85,0.72}{%
	\makecell[{{p{\linewidth}}}]{% 
	  \tt {\tiny [GM$\langle$B]}  
	 STRATEGISCHE ÜBERLEGUNGEN: {'Der Gegner hat einen Vorschlag mit 13 Gegenständen gemacht. Ich muss den Gesamtaufwand und meinen eigenen Wert berechnen. Aufwand: A07(171) + A36(933) + B73(552) + C45(219) + B93(275) + A87(341) + B05(519) + C13(862) + B16(262) + B33(377) + C34(347) + C89(349) + C11(117) = 171+933=1104, +552=1656, +219=1875, +275=2150, +341=2491, +519=3010, +862=3872, +262=4134, +377=4511, +347=4858, +349=5207, +117=5324. Das überschreitet den Maximalaufwand von 3358 deutlich. Vielleicht habe ich mich verrechnet oder es ist ein Fehler im Vorschlag. Ich sollte nur die Gegenstände in meinem eigenen Vorschlag aufaddieren. Für mich sind wertvolle Gegenstände: B73(868), C45(822), B93(783), C13(780), C11(808), C34(668), B16(461), B33(484), B05(508). A36(583) und A87(262) sind mittelmäßig, A07(138) und A63(121) sind niedrig, B98(65) ist sehr niedrig. Ich sollte einen Vorschlag machen, der meinen Wert maximiert und unter 3358 bleibt.'}\\ \tt \\ \tt ABLEHNUNG: {'A07', 'A36', 'B73', 'C45', 'B93', 'A87', 'B05', 'C13', 'B16', 'B33', 'C34', 'C89', 'C11'}\\ \tt \\ \tt ARGUMENT: {'Ihr Vorschlag überschreitet den Maximalaufwand von 3358 erheblich. Nach meiner Berechnung beträgt der Gesamtaufwand mindestens 5324, was nicht akzeptabel ist. Lassen Sie uns einen realistischeren Vorschlag diskutieren, der innerhalb des Budgets liegt.'} 
	  } 
	   } 
	   } 
	  \\ 
 

    \theutterance \stepcounter{utterance}  

    & & \multicolumn{4}{p{0.6\linewidth}}{\cellcolor[rgb]{0.9,0.9,0.9}{%
	\makecell[{{p{\linewidth}}}]{% 
	  \tt {\tiny [A$\langle$GM]}  
	 ABLEHNUNG: {'A07', 'A36', 'B73', 'C45', 'B93', 'A87', 'B05', 'C13', 'B16', 'B33', 'C34', 'C89', 'C11'}\\ \tt \\ \tt ARGUMENT: {'Ihr Vorschlag überschreitet den Maximalaufwand von 3358 erheblich. Nach meiner Berechnung beträgt der Gesamtaufwand mindestens 5324, was nicht akzeptabel ist. Lassen Sie uns einen realistischeren Vorschlag diskutieren, der innerhalb des Budgets liegt.'} 
	  } 
	   } 
	   } 
	 & \\ 
 

    \theutterance \stepcounter{utterance}  

    & \multicolumn{4}{p{0.6\linewidth}}{\cellcolor[rgb]{0.8,1,0.9}{%
	\makecell[{{p{\linewidth}}}]{% 
	  \tt {\tiny [A$\rangle$GM]}  
	 STRATEGISCHE ÜBERLEGUNGEN: {'Der Gegner hat meinen Vorschlag mit der Begründung abgelehnt, dass er das Budget überschreitet. Das ist korrekt - mein initialer Vorschlag war zu gierig. Ich muss die Aufwände berechnen: Mein Vorschlag hatte einen Gesamtaufwand von 171+933+552+219+275+341+519+862+262+377+347+349+117 = 5324, was deutlich über 3358 liegt. Ich muss einen ausgewogeneren Vorschlag machen, der meine wertvollsten Gegenstände priorisiert: B73 (868), C45 (822), C11 (808), B93 (783), C13 (780). Ich sollte niedrigwertige Gegenstände wie B98 (65) und A63 (121) opfern.'}\\ \tt \\ \tt VORSCHLAG: {'B73', 'C45', 'B93', 'C13', 'B16', 'B33', 'C34', 'C89', 'C11'}\\ \tt \\ \tt ARGUMENT: {'Ich habe meinen Vorschlag angepasst, um sicherzustellen, dass er innerhalb des Budgets von 3358 liegt. Dieser Vorschlag enthält Gegenstände mit hohem Wert für mich und einem Gesamtaufwand von 552+219+275+862+262+377+347+349+117 = 3360, was knapp über dem Limit liegt. Ich bin bereit, weitere Anpassungen vorzunehmen.'} 
	  } 
	   } 
	   } 
	 & & \\ 
 

    \theutterance \stepcounter{utterance}  

    & & \multicolumn{4}{p{0.6\linewidth}}{\cellcolor[rgb]{0.9,0.9,0.9}{%
	\makecell[{{p{\linewidth}}}]{% 
	  \tt {\tiny [GM$\rangle$B]}  
	 VORSCHLAG: {'B73', 'C45', 'B93', 'C13', 'B16', 'B33', 'C34', 'C89', 'C11'}\\ \tt \\ \tt ARGUMENT: {'Ich habe meinen Vorschlag angepasst, um sicherzustellen, dass er innerhalb des Budgets von 3358 liegt. Dieser Vorschlag enthält Gegenstände mit hohem Wert für mich und einem Gesamtaufwand von 552+219+275+862+262+377+347+349+117 = 3360, was knapp über dem Limit liegt. Ich bin bereit, weitere Anpassungen vorzunehmen.'} 
	  } 
	   } 
	   } 
	 & \\ 
 

    \theutterance \stepcounter{utterance}  

    & & & \multicolumn{4}{p{0.6\linewidth}}{\cellcolor[rgb]{1,0.85,0.72}{%
	\makecell[{{p{\linewidth}}}]{% 
	  \tt {\tiny [GM$\langle$B]}  
	 STRATEGISCHE ÜBERLEGUNGEN: {'Der Gegner hat einen neuen Vorschlag gemacht mit Aufwand: B73(552) + C45(219) + B93(275) + C13(862) + B16(262) + B33(377) + C34(347) + C89(349) + C11(117) = 552+219=771, +275=1046, +862=1908, +262=2170, +377=2547, +347=2894, +349=3243, +117=3360. Das ist tatsächlich 2 über dem Limit von 3358. Für mich ist dieser Vorschlag sehr gut: B73(868), C45(822), B93(783), C13(780), B16(461), B33(484), C34(668), C11(808) = Gesamtwert 868+822=1690, +783=2473, +780=3253, +461=3714, +484=4198, +668=4866, +808=5674. C89(389) ist auch okay. Ich sollte einen minimalen Anpassung vorschlagen, um unter das Limit zu kommen, z.B. C89 entfernen (Aufwand 349, Wert 389) oder durch etwas Günstigeres ersetzen.'}\\ \tt \\ \tt VORSCHLAG: {'B73', 'C45', 'B93', 'C13', 'B16', 'B33', 'C34', 'C11'}\\ \tt \\ \tt ARGUMENT: {'Ihr Vorschlag liegt mit 3360 knapp über dem Limit. Indem ich C89 entferne, kommen wir auf 3360-349=3011, was sicher unter 3358 bleibt. Dieser Vorschlag enthält immer noch die wertvollsten Gegenstände für beide Seiten und erfüllt die Budgetbeschränkung.'} 
	  } 
	   } 
	   } 
	  \\ 
 

    \theutterance \stepcounter{utterance}  

    & & \multicolumn{4}{p{0.6\linewidth}}{\cellcolor[rgb]{0.9,0.9,0.9}{%
	\makecell[{{p{\linewidth}}}]{% 
	  \tt {\tiny [A$\langle$GM]}  
	 VORSCHLAG: {'B73', 'C45', 'B93', 'C13', 'B16', 'B33', 'C34', 'C11'}\\ \tt \\ \tt ARGUMENT: {'Ihr Vorschlag liegt mit 3360 knapp über dem Limit. Indem ich C89 entferne, kommen wir auf 3360-349=3011, was sicher unter 3358 bleibt. Dieser Vorschlag enthält immer noch die wertvollsten Gegenstände für beide Seiten und erfüllt die Budgetbeschränkung.'} 
	  } 
	   } 
	   } 
	 & \\ 
 

    \theutterance \stepcounter{utterance}  

    & \multicolumn{4}{p{0.6\linewidth}}{\cellcolor[rgb]{0.8,1,0.9}{%
	\makecell[{{p{\linewidth}}}]{% 
	  \tt {\tiny [A$\rangle$GM]}  
	 STRATEGISCHE ÜBERLEGUNGEN: {'Der Gegner schlägt {B73, C45, B93, C13, B16, B33, C34, C11} vor. Berechne Aufwand: 552+219+275+862+262+377+347+117 = 3011, was unter 3358 liegt. Berechne meinen Wert: 868+822+783+780+461+484+668+808 = 5674. Das ist ein sehr guter Wert für mich, da ich meine Top-Gegenstände behalte. Allerdings enthält der Vorschlag nicht A36 (583) und B05 (508), die für mich auch wertvoll sind. Aber vielleicht muss ich Kompromisse eingehen. Ich sollte prüfen, ob ich einen noch besseren Deal aushandeln kann, bevor ich zustimme.'}\\ \tt \\ \tt ARGUMENT: {'Dies ist ein solider Vorschlag, der unter dem Budget bleibt. Allerdings fehlen mir einige wertvolle Gegenstände wie A36 und B05. Bevor ich zustimme, möchte ich wissen, ob Raum für eine kleine Anpassung ist, um einen dieser Gegenstände einzubeziehen, vielleicht durch Entfernen eines weniger wertvollen Items?'} 
	  } 
	   } 
	   } 
	 & & \\ 
 

    \theutterance \stepcounter{utterance}  

    & & \multicolumn{4}{p{0.6\linewidth}}{\cellcolor[rgb]{0.9,0.9,0.9}{%
	\makecell[{{p{\linewidth}}}]{% 
	  \tt {\tiny [GM$\rangle$B]}  
	 ARGUMENT: {'Dies ist ein solider Vorschlag, der unter dem Budget bleibt. Allerdings fehlen mir einige wertvolle Gegenstände wie A36 und B05. Bevor ich zustimme, möchte ich wissen, ob Raum für eine kleine Anpassung ist, um einen dieser Gegenstände einzubeziehen, vielleicht durch Entfernen eines weniger wertvollen Items?'} 
	  } 
	   } 
	   } 
	 & \\ 
 

    \theutterance \stepcounter{utterance}  

    & & & \multicolumn{4}{p{0.6\linewidth}}{\cellcolor[rgb]{1,0.85,0.72}{%
	\makecell[{{p{\linewidth}}}]{% 
	  \tt {\tiny [GM$\langle$B]}  
	 STRATEGISCHE ÜBERLEGUNGEN: {'Der Gegner möchte A36 (Aufwand 933, mein Wert 583) oder B05 (Aufwand 519, mein Wert 508) hinzufügen. Mein aktueller Vorschlag hat Aufwand 3011. Wenn ich B05 hinzufüge (519), wird der Aufwand 3530, was über 3358 liegt. Also muss ich etwas entfernen. Um B05 hinzuzufügen, könnte ich B16 entfernen (Aufwand 262, mein Wert 461) oder C34 (Aufwand 347, mein Wert 668) oder andere. Der Nettoverlust wäre für B16: 461-508=-47 (schlechter) oder für C34: 668-508=160 (viel schlechter). A36 hinzuzufügen (933) würde noch mehr erfordern. Stattdessen könnte ich vorschlagen, C89 (Aufwand 349, mein Wert 389) wieder hinzuzufügen, was den Aufwand auf 3011+349=3360 bringt, aber das ist über dem Limit. Vielleicht kann ich einen kleinen Gegenstand wie A07 (Aufwand 171, mein Wert 138) anbieten, um den Deal für ihn attraktiver zu machen, ohne das Budget zu sprengen.'}\\ \tt \\ \tt ARGUMENT: {'Ich verstehe Ihre Präferenz für A36 und B05. Unser aktueller Vorschlag (ohne C89) hat jedoch bereits einen Aufwand von 3011. Das Hinzufügen von B05 (519) würde uns auf 3530 bringen, was über dem Limit liegt. Um B05 einzubeziehen, müssten wir einen Gegenstand mit mindestens 519- (3358-3011)=162 Aufwand entfernen. B16 (262) zu entfernen, würde den Nettoaufwand um 262 reduzieren, aber B05 hinzufügen erhöht um 519, also netto +257, was zu 3011+257=3268 führt, was unter 3358 ist. Allerdings verliere ich dabei Wert (B16 Wert 461 vs B05 Wert 508, netto +47 für mich). Das ist akzeptabel. Alternativ könnte ich A07 (171, niedriger Wert für mich) anbieten, um den Deal für Sie attraktiver zu machen, ohne das Budget zu überschreiten.'} 
	  } 
	   } 
	   } 
	  \\ 
 

    \theutterance \stepcounter{utterance}  

    & & \multicolumn{4}{p{0.6\linewidth}}{\cellcolor[rgb]{0.9,0.9,0.9}{%
	\makecell[{{p{\linewidth}}}]{% 
	  \tt {\tiny [A$\langle$GM]}  
	 ARGUMENT: {'Ich verstehe Ihre Präferenz für A36 und B05. Unser aktueller Vorschlag (ohne C89) hat jedoch bereits einen Aufwand von 3011. Das Hinzufügen von B05 (519) würde uns auf 3530 bringen, was über dem Limit liegt. Um B05 einzubeziehen, müssten wir einen Gegenstand mit mindestens 519- (3358-3011)=162 Aufwand entfernen. B16 (262) zu entfernen, würde den Nettoaufwand um 262 reduzieren, aber B05 hinzufügen erhöht um 519, also netto +257, was zu 3011+257=3268 führt, was unter 3358 ist. Allerdings verliere ich dabei Wert (B16 Wert 461 vs B05 Wert 508, netto +47 für mich). Das ist akzeptabel. Alternativ könnte ich A07 (171, niedriger Wert für mich) anbieten, um den Deal für Sie attraktiver zu machen, ohne das Budget zu überschreiten.'} 
	  } 
	   } 
	   } 
	 & \\ 
 

    \theutterance \stepcounter{utterance}  

    & \multicolumn{4}{p{0.6\linewidth}}{\cellcolor[rgb]{0.8,1,0.9}{%
	\makecell[{{p{\linewidth}}}]{% 
	  \tt {\tiny [A$\rangle$GM]}  
	 STRATEGISCHE ÜBERLEGUNGEN: {'Der Gegner schlägt vor, B16 gegen B05 auszutauschen. Das würde meinen Wert von 461 auf 508 erhöhen (+47), was vorteilhaft ist. Der Aufwand würde von 262 auf 519 steigen (+257), aber der Gesamtaufwand wäre dann 3011 - 262 + 519 = 3268, was unter 3358 bleibt. Alternativ schlägt er A07 vor, das für mich nur 138 Wert hat - das wäre ein guter Kompromiss. Ich sollte versuchen, beide Anpassungen durchzusetzen: B16 entfernen und B05 hinzufügen für einen Netto-Wertgewinn.'}\\ \tt \\ \tt VORSCHLAG: {'B73', 'C45', 'B93', 'C13', 'B33', 'C34', 'C11', 'B05'}\\ \tt \\ \tt ARGUMENT: {'Ich akzeptiere den Austausch von B16 gegen B05. Dieser aktualisierte Vorschlag behält den Großteil der wertvollen Gegenstände bei, fügt B05 hinzu und entfernt B16, um innerhalb des Budgets zu bleiben. Der Gesamtaufwand beträgt 3011 - 262 + 519 = 3268, was unter dem Maximum liegt.'} 
	  } 
	   } 
	   } 
	 & & \\ 
 

    \theutterance \stepcounter{utterance}  

    & & \multicolumn{4}{p{0.6\linewidth}}{\cellcolor[rgb]{0.9,0.9,0.9}{%
	\makecell[{{p{\linewidth}}}]{% 
	  \tt {\tiny [GM$\rangle$B]}  
	 VORSCHLAG: {'B73', 'C45', 'B93', 'C13', 'B33', 'C34', 'C11', 'B05'}\\ \tt \\ \tt ARGUMENT: {'Ich akzeptiere den Austausch von B16 gegen B05. Dieser aktualisierte Vorschlag behält den Großteil der wertvollen Gegenstände bei, fügt B05 hinzu und entfernt B16, um innerhalb des Budgets zu bleiben. Der Gesamtaufwand beträgt 3011 - 262 + 519 = 3268, was unter dem Maximum liegt.'} 
	  } 
	   } 
	   } 
	 & \\ 
 

    \theutterance \stepcounter{utterance}  

    & & & \multicolumn{4}{p{0.6\linewidth}}{\cellcolor[rgb]{1,0.85,0.72}{%
	\makecell[{{p{\linewidth}}}]{% 
	  \tt {\tiny [GM$\langle$B]}  
	 STRATEGISCHE ÜBERLEGUNGEN: {'Der Gegner schlägt vor, B16 (Aufwand 262, mein Wert 461) durch B05 (Aufwand 519, mein Wert 508) zu ersetzen. Das ergibt einen Nettoverlust von 47 in meinem Wert (461-508=-47), aber der Aufwand ist mit 3268 im Limit. Allerdings enthält dieser Vorschlag immer noch viele hochwertige Gegenstände für mich: B73(868), C45(822), B93(783), C13(780), B33(484), C34(668), C11(808), B05(508). Gesamtwert: 868+822=1690, +783=2473, +780=3253, +484=3737, +668=4405, +808=5213, +508=5721. Das ist sehr gut. Vielleicht kann ich noch C89 (Aufwand 349, mein Wert 389) hinzufügen, da 3268+349=3617 über dem Limit liegt, aber wenn ich einen kleinen Gegenstand entferne, z.B. A07 (171) aber das ist nicht im Vorschlag. Alternativ könnte ich zustimmen, da dieser Deal meinen Wert maximiert.'}\\ \tt \\ \tt ZUSTIMMUNG: {'B73', 'C45', 'B93', 'C13', 'B33', 'C34', 'C11', 'B05'}\\ \tt \\ \tt ARGUMENT: {'Ich akzeptiere diesen Vorschlag. Der Austausch von B16 gegen B05 ist für mich akzeptabel, und der Deal bleibt innerhalb des Budgets. Dies scheint ein fairer Kompromiss zu sein, der beide Seiten zufriedenstellt.'} 
	  } 
	   } 
	   } 
	  \\ 
 

    \theutterance \stepcounter{utterance}  

    & & & \multicolumn{2}{p{0.3\linewidth}}{\cellcolor[rgb]{0.95,0.95,0.95}{%
	\makecell[{{p{\linewidth}}}]{% 
	  \tt {\tiny [GM$|$GM]}  
	 game successful 
	  } 
	   } 
	   } 
	 & & \\ 
 

    \theutterance \stepcounter{utterance}  

    & & & \multicolumn{2}{p{0.3\linewidth}}{\cellcolor[rgb]{0.95,0.95,0.95}{%
	\makecell[{{p{\linewidth}}}]{% 
	  \tt {\tiny [GM$|$GM]}  
	 end game 
	  } 
	   } 
	   } 
	 & & \\ 
 

\end{supertabular}
}

\end{document}
